\section{The marking and the exact criterion}
Let $I=\{0,\ldots,r-1\}$, where $r\geq3$, and let
\[
 G=\langle\alpha_1,\ldots,\alpha_{r-1}\rangle,
 \qquad\alpha_0=1,\qquad a_{ij}=\alpha_i^{-1}\alpha_j,
 \qquad P_S=\langle a_{ij}:i,j\in S\rangle.
\]
We use the $C^+$ marking of Zhang--Leemans: the specified nonidentity generators are independent, and
\begin{equation}\label{eq:IC}
                 P_S\cap P_T=P_{S\cap T}
                  \qquad(|S|,|T|\geq2).
\end{equation}
Singleton and empty supports generate the trivial group. In particular the $\alpha_i$ are distinct. Put $H_i=P_{I\setminus\{i\}}$ and let $\Gamma$ be their coset incidence geometry. The word geometry is a hypothesis: every flag extends to a chamber. Our left cosets correspond to the source's right cosets by $gH_i\mapsto H_i g^{-1}$; this bijection preserves nonempty intersection and converts the acting element to its inverse.

Repeated use of \eqref{eq:IC} gives
\begin{equation}\label{eq:parabolics}
              H_J:=\bigcap_{j\in J}H_j=P_{I\setminus J}.
\end{equation}
Indeed, intersect successively while the remaining support has at least two elements; once its size is one, the group is already trivial. Thus $H_J=1$ whenever $|J|\geq r-1$. The panel-stabilizer condition listed separately in source Theorem 1.1 is consequently automatic under this marking hypothesis.

\begin{theorem}\label{thm:criterion}
The geometry $\Gamma$ is a chiral hypertope if and only if the following conditions hold:
\begin{enumerate}[label=(\roman*)]
\item for some $k\in I$, $G$ is transitive on the chambers of the truncation omitting type $k$;
\item for that $k$, the normalized extension set is exactly two cosets:
\begin{equation}\label{eq:twocosets}
 \bigcap_{j\ne k}H_jH_k=H_k\ \dot\cup\ xH_k
                         \quad\text{with }x\notin H_k;
\end{equation}
\item no automorphism of $G$ sends every $\alpha_i$ to $\alpha_i^{-1}$.
\end{enumerate}
When these conditions hold, condition (i) and the corresponding two-coset condition hold for every $k$. If the groups are finite, (ii) is equivalent to the product-size condition in Zhang--Leemans, Theorem 1.1. For infinite groups the exact coset condition, not that cardinal equation, is the statement proved here.
\end{theorem}

\section{Neighbors and proper flags}
Let $C=\{H_i:i\in I\}$ be the base chamber. For $i\in I$ define
\begin{equation}\label{eq:neighbor}
 D_i=(C\setminus\{H_i\})\cup\{a_{ij}H_i\},\qquad j\ne i.
\end{equation}
The chosen $j$ does not matter: $a_{ij}^{-1}a_{ik}=a_{jk}\in H_i$ when $j,k\ne i$. This is a chamber. To check its modified object against $H_l$, choose a representative $a_{ij}$ with $j\ne i,l$; this is possible because $r\geq3$, and the representative lies in $H_l$. Also $a_{ij}\notin H_i$: otherwise \eqref{eq:IC} applied to $\{i,j\}$ and $I\setminus\{i\}$ would put $a_{ij}$ in $P_{\{j\}}=1$, contrary to the distinctness of $\alpha_i,\alpha_j$. Thus $C$ and $D_i$ are distinct and differ only at type $i$.

The key transport identity is
\begin{equation}\label{eq:neighbortransport}
                         a_{ji}D_i=D_j.
\end{equation}
At type $i$ the multiplier cancels $a_{ij}$; at type $j$ it creates the modified coset of $D_j$. At every other type $l$, both $i,j$ differ from $l$, so $a_{ji}\in H_l$ and the coset remains $H_l$. This verifies the identity at every coordinate.

Assume (i)--(ii). The extension-fiber theorem in the common paper, together with \eqref{eq:parabolics}, gives exactly two free chamber orbits. They are represented by $C$ and $D_k$: no element of $G$ can carry $C$ to $D_k$, because it would fix their panel, whose stabilizer is trivial. Formula~\eqref{eq:neighbortransport} puts all $D_i$ in the second orbit.

Every proper flag is in the orbit of a flag of $C$. To see this, extend it to a chamber and move that chamber to $C$ or $D_k$. In the second case choose a type $i$ missing from the flag and move $D_k$ to $D_i$ by \eqref{eq:neighbortransport}. Its objects of the flag's types then coincide with those of $C$. Hence $G$ is transitive on every proper flag type.

Fix now the panel of $C$ missing $i$. It has the two completions $C,D_i$. Any other completion is in one of their two free orbits. An element carrying the appropriate representative to that completion fixes the whole panel and must be the identity. The panel therefore has exactly those two completions. Transitivity on panels proves thinness everywhere. This argument counts completions within each orbit and is valid for infinite $G$; it does not infer a panel size by dividing infinite cardinalities.

\section{Every residue is connected}
Normalize a flag of type $J$ to $\{H_j:j\in J\}$ and set $S=I\setminus J$. For $i\in S$, transitivity on flags of type $J\cup\{i\}$ shows that the type-$i$ objects of the residue are exactly
\[
                  H_J/(H_J\cap H_i).
\]
If $|S|\geq3$, this assertion uses only proper flag transitivity. The same is true when $|S|=2$, because $|J\cup\{i\}|=r-1$. The internal coset incidence system has all these objects and some of the residue's edges; it is a spanning subgraph, not automatically an incidence isomorphism.

When $|S|\geq3$, its defining subgroups are $P_{S\setminus\{i\}}$. They generate $P_S=H_J$, since each generator $a_{uv}$ belongs to $P_{S\setminus\{i\}}$ for an $i$ different from $u,v$. For completeness, generation implies connectivity as follows. The objects represented by the identity are mutually incident. Multiplying their whole chamber by an element of one defining subgroup leaves that subgroup's object fixed, so the two chambers lie in the same connected component. Induction on a word in the defining subgroups brings every group translate into this component. These translates contain every object. The larger actual residue graph is therefore connected too.

Suppose $S=\{u,v\}$. Then $H_J=\langle a_{uv}\rangle$ is cyclic and both vertex stabilizers are trivial. The residue has two copies of this cyclic group. It contains the matching edges
\[
        h_{\rm left}\,--\,h_{\rm right},\qquad
        h_{\rm left}\,--\,(h a_{uv}^{-1})_{\rm right}.
\]
The first edges come from $hC$. The second come from $hD_v$, since $D_v$ has its type-$v$ object represented by $a_{vu}=a_{uv}^{-1}$. Two successive matching steps move the cyclic coordinate by $a_{uv}$ or its inverse, so these edges connect all vertices. They form a cycle in the finite case and a doubly infinite path in the infinite case. Any additional incidence edges preserve connectivity. This proves residual connectedness for every flag and completes the hypertope part of the sufficient direction.

\section{The full automorphism group and the converse}
Take $D_0$ as the negative base chamber. Equation~\eqref{eq:neighbortransport} gives
\[
                         D_i=\alpha_i^{-1}D_0.
\]
Label the two free chamber orbits by $(g,+)=gC$ and $(g,-)=gD_0$. Thinness gives exactly the adjacency formula
\[
                   (g,+)\sim_i(g\alpha_i^{-1},-)
\]
used by the common reflection theorem. That theorem proves that (iii) is equivalent to chirality and identifies the full type-preserving automorphism group, rather than assuming it equals the designated group $G$.

For necessity, suppose $\Gamma$ is a chiral hypertope, with chirality defined using its full type-preserving automorphism group. We first show directly that the designated $G$ has the two chamber orbits required above. The chambers $D_i$ are the unique $i$-neighbors of $C$ by thinness. Let $\mathcal C$ be the union of the $G$-orbits of $C$ and $D_0$. Every neighbor of $gC$ is $gD_i$, in the second orbit. Conversely every $i$-neighbor of $gD_0$ is $g\alpha_i C$: translate the $i$-edge from $C$ to $D_i=\alpha_i^{-1}D_0$ by $\alpha_i$. Thus $\mathcal C$ is closed under adjacency. Residual connectedness implies chamber connectedness by the common propagation lemma, so $\mathcal C$ is the entire chamber set.

The two orbits are distinct. If $gC=D_0$, then $g$ fixes their common panel and lies in $H_{I\setminus\{0\}}=1$, which is impossible. The stabilizer of a chamber is also trivial. The proper-flag argument of Section 2 now gives transitivity on every truncation. Normalizing a panel, its exactly two completions give precisely two cosets by the extension-fiber theorem, proving (i)--(ii) for every $k$. Finally the common reflection theorem would make the full geometry regular if (iii) failed, contrary to chirality. This proves necessity without presupposing that the designated action is the full chiral action.

\section{Source scope and attribution}
The exact principal target is Theorem 1.1 of W.-J. Zhang and D. Leemans, \emph{Testing chirality on hypertopes}, arXiv:2511.14562v1, the ten-page November 2025 manuscript. All ten acquired pages have been read for this edition. The source states the cardinality test without an explicit finiteness restriction; the two-coset formulation above records the necessary infinite-group qualification. Condition (ii) concerning panel stabilizers in the source is retained as a proved consequence of its $C^+$ hypothesis.

The orbit normalization also explains the source's Theorem 3.2, but this edition does not enlarge the admitted principal target to a replacement of every ancillary result or the three Magma functions printed in Section 5. Those functions and their software dependencies are not required by the uniform written proof here. The source's passage from a marked inversion to a geometric reflection cites Fernandes--Leemans--Weiss (2016), Theorem 8.2; the common paper supplies the actual chamber map and its geometric reconstruction. The residue argument reuses the earlier reviewed proof bank's distinction between an internal coset graph and the larger actual residue, including the second matching in rank two. No finite experiment is used to establish the quantified theorem, and no historical novelty is asserted for the criterion.
